% !TEX program = xelatex
\documentclass[10pt,letterpaper,twocolumn]{article}

% === GEOMETRY ===
\usepackage[top=0.75in, bottom=0.75in, left=0.75in, right=0.75in]{geometry}

% === FONTS ===
\usepackage{fontspec}
\usepackage{unicode-math}

% Latin Modern has NO Cyrillic glyphs. STIX Two Text is a Latin-Modern-
% adjacent serif with full Cyrillic that ships on macOS AND in TeX Live
% (stix2-otf), so this preamble renders identically on macOS xelatex and
% the oven's Linux xelatex.
\setmainfont{STIX Two Text}
\setsansfont{STIX Two Text}

% Custom AC fonts
\newfontfamily\acbold{ywft-processing-bold}[
  Path=../../system/public/type/webfonts/,
  Extension=.ttf
]
\newfontfamily\aclight{ywft-processing-light}[
  Path=../../system/public/type/webfonts/,
  Extension=.ttf
]
\setmonofont{Latin Modern Mono}[
  Extension=.otf,
  UprightFont=lmmono10-regular,
  ItalicFont=lmmono10-italic,
  Scale=0.85,
]

% === PACKAGES ===
\usepackage{xcolor}
\usepackage{titlesec}
\usepackage{enumitem}
\usepackage{booktabs}
\usepackage{tabularx}
\usepackage{multicol}
\usepackage{fancyhdr}
\usepackage{hyperref}
\usepackage{xurl}
\usepackage{graphicx}
\graphicspath{{figures/}{../../papers/arxiv-ac/figures/}}
\usepackage{ragged2e}
\usepackage{microtype}
\usepackage{listings}
\usepackage{natbib}
\usepackage[colorspec=0.92]{draftwatermark}

% === COLORS (AC palette) ===
\definecolor{acpink}{RGB}{180,72,135}
\definecolor{acpurple}{RGB}{120,80,180}
\definecolor{acdark}{RGB}{64,56,74}
\definecolor{acgray}{RGB}{119,119,119}
\definecolor{draftcolor}{RGB}{180,72,135}

% === DRAFT WATERMARK ===
\DraftwatermarkOptions{
  text=WORKING DRAFT,
  fontsize=3cm,
  color=draftcolor!18,
  angle=45,
  pos={0.5\paperwidth, 0.5\paperheight}
}

% === JS SYNTAX COLORS ===
\definecolor{jskw}{RGB}{119,51,170}
\definecolor{jsfn}{RGB}{0,136,170}
\definecolor{jsstr}{RGB}{170,120,0}
\definecolor{jsnum}{RGB}{204,0,102}
\definecolor{jscmt}{RGB}{102,102,102}

% === HYPERREF ===
\hypersetup{
  colorlinks=true,
  linkcolor=acpurple,
  urlcolor=acpurple,
  citecolor=acpurple,
  pdfauthor={@jeffrey},
  pdftitle={PLOrk'ing the Planet: От оркестра ноутбуков к планетарному инструменту},
}

% === SECTION FORMATTING ===
\titleformat{\section}
  {\normalfont\bfseries\normalsize\uppercase}
  {\thesection.}
  {0.5em}
  {}
\titlespacing{\section}{0pt}{1.2em}{0.3em}

\titleformat{\subsection}
  {\normalfont\bfseries\small}
  {\thesubsection}
  {0.5em}
  {}
\titlespacing{\subsection}{0pt}{0.8em}{0.2em}

% === HEADER/FOOTER ===
\pagestyle{fancy}
\fancyhf{}
\renewcommand{\headrulewidth}{0pt}
\fancyhead[C]{\footnotesize\color{acpink}\textit{Рабочий черновик --- не цитировать}}
\fancyfoot[C]{\footnotesize\thepage}

% === CUSTOM COMMANDS ===
\newcommand{\acdot}{{\color{acpink}.}}
\newcommand{\ac}{\textsc{Aesthetic.Computer}}
% Random caps for Aesthetic.Computer branding
\newcount\acrandtmp
\newcommand{\acrandletter}[2]{%
  \acrandtmp=\uniformdeviate 2\relax
  \ifnum\acrandtmp=0\relax#1\else#2\fi%
}
\newcommand{\acrandname}{%
  \acrandletter{a}{A}\acrandletter{e}{E}\acrandletter{s}{S}\acrandletter{t}{T}%
  \acrandletter{h}{H}\acrandletter{e}{E}\acrandletter{t}{T}\acrandletter{i}{I}%
  \acrandletter{c}{C}{\color{acpink}.}\acrandletter{c}{C}\acrandletter{o}{O}%
  \acrandletter{m}{M}\acrandletter{p}{P}\acrandletter{u}{U}\acrandletter{t}{T}%
  \acrandletter{e}{E}\acrandletter{r}{R}%
}
\newcommand{\acos}{\textsc{AC Native OS}}

% === LISTINGS ===
\lstdefinelanguage{acjs}{
  morekeywords=[1]{function,export,const,let,var,return,if,else,new,async,await,import,from},
  morekeywords=[2]{wipe,ink,line,box,circle,write,screen,params,colon,jump,send,store,net,sound,speaker,system},
  sensitive=true,
  morecomment=[l]{//},
  morestring=[b]",
  morestring=[b]',
  morestring=[b]`,
  escapeinside={|}{|},
}

\lstdefinestyle{acjsstyle}{
  language=acjs,
  keywordstyle=[1]\color{jskw}\bfseries,
  keywordstyle=[2]\color{jsfn}\bfseries,
  commentstyle=\color{jscmt}\itshape,
  stringstyle=\color{jsstr},
}

\lstset{
  basicstyle=\ttfamily\small,
  breaklines=true,
  frame=single,
  rulecolor=\color{acgray!30},
  backgroundcolor=\color{acgray!5},
  xleftmargin=0.5em,
  xrightmargin=0.5em,
  aboveskip=0.5em,
  belowskip=0.5em,
}

% === LIST SETTINGS ===
\setlist[itemize]{nosep, leftmargin=1.2em, itemsep=0.1em}
\setlist[enumerate]{nosep, leftmargin=1.2em}

% === COLUMN SEPARATION ===
\setlength{\columnsep}{1.8em}

% === PARAGRAPH SETTINGS ===
\setlength{\parindent}{1em}
\setlength{\parskip}{0.3em}

% Narrow two-column layout. Russian is set without hyphenation patterns
% here, so long words are unbreakable — give the line-breaker much more
% slack to avoid words spilling into the column gutter.
\tolerance=4000
\emergencystretch=4em
\hyphenpenalty=50
\hbadness=10000

\begin{document}

% ============ TITLE BLOCK ============

\twocolumn[{%
\begin{center}

\includegraphics[height=4em]{pals}\par\vspace{0.5em}
{\acbold\fontsize{24pt}{28pt}\selectfont\color{acdark} PLOrk'ing the Planet}\par
\includegraphics[width=0.5\textwidth]{figures/cover}\par
\vspace{0.35em}
\vspace{0.2em}
% AC display fonts (ywft) are Latin-only — use the Cyrillic main font
% for the Russian subtitle/tagline so they don't render as tofu.
{\rmfamily\itshape\fontsize{12pt}{15pt}\selectfont\color{acpink} От оркестра ноутбуков к планетарному инструменту}\par
\vspace{0.3em}
{\rmfamily\fontsize{9.5pt}{12pt}\selectfont\color{acgray} Списанное железо, синтез на «голом железе» и оркестр из 240 миллионов машин}\par
\vspace{0.6em}
{\normalsize\href{https://prompt.ac/@jeffrey}{@jeffrey}}\par
{\small\color{acgray} Aesthetic.Computer}\par
{\small\color{acgray} ORCID: \href{https://orcid.org/0009-0007-4460-4913}{0009-0007-4460-4913}}\par
\vspace{0.3em}
{\small\color{acpurple} \url{https://aesthetic.computer}}\par
\vspace{0.6em}
\rule{\textwidth}{1.5pt}
\vspace{0.5em}
\end{center}

\begin{center}
{\small\color{acpink}\textbf{[ рабочий черновик --- не цитировать ]}}
\end{center}
\vspace{0.3em}

\begin{quote}
\small\noindent\textbf{Аннотация.}
В 2005 году Дэн Труман и Перри Кук основали Princeton Laptop Orchestra (PLOrk) --- пятнадцать станций, каждая из которых состояла из MacBook, заказного полусферического громкоговорителя, 8-канальной стойки усилителя и набора контроллеров, работавших на ChucK и Max/MSP. PLOrk доказал, что ансамбли ноутбуков --- музыкально полноценная форма исполнения, и породил более тридцати университетских оркестров ноутбуков по всему миру. Но двадцать лет спустя каждый из этих ансамблей по-прежнему заперт внутри университета: ограничен зачислением, финансируется кафедрами, размещён в специальных студиях и обслуживается техническим персоналом. Оркестр ноутбуков так и не вышел за пределы кампуса. Мы утверждаем, что модель PLOrk содержит невозможность масштабирования: каждое место стоит \$1.500+, требует изготовления заказных громкоговорителей, зависит от институциональной инфраструктуры и обслуживает максимум 15 человек в семестр. L2Ork (Virginia Tech) снизил стоимость до \$750 за место с помощью Linux-нетбуков и более дешёвых полусферических громкоговорителей, но остался привязан к университету. Мечта об оркестре ноутбуков --- что сетевые компьютеры могут стать новым видом ансамблевого инструмента --- всегда была мечтой о доступе. Она провалилась именно в доступе. Эта статья представляет конкретную альтернативу. \acos{} --- операционная система объёмом 89\,МБ, которая загружает любой списанный ноутбук x86\_64 в готовый музыкальный инструмент за 7,3 секунды: встроенный 32-голосный полифонический синтезатор, сетевое ансамблевое соединение «из коробки», персональная идентичность и обновления по воздуху. Списанный ThinkPad стоит \$50. Запись на USB занимает минуту. Без полусферических громкоговорителей, без стоек усилителей, без лицензий Max/MSP, без ИТ-отдела. С 240 миллионами ПК, объявленными устаревшими из-за окончания поддержки Windows~10, и 62 миллионами тонн электронных отходов, образующихся ежегодно, сырьё для планетарного оркестра ноутбуков уже существует. Вопрос не в том, можем ли мы позволить себе «PLOrk'ить планету». Вопрос в том, можем ли мы позволить себе этого не делать.
\end{quote}
\vspace{0.5em}
}]

% ============ 1. ВВЕДЕНИЕ ============

\section{Введение}

Princeton Laptop Orchestra не была группой с ноутбуками. Это было утверждение: что ноутбук может быть оркестровым инструментом, что программное обеспечение может производить пространственно богатый звук, что сеть компьютеров может составлять ансамбль столь же полноценный, как струнный квартет. Дэн Труман, Перри Кук, Скотт Смолвуд и Гэ Ван построили PLOrk в 2005 году в Принстонском университете~\citep{trueman2006plork, trueman2007why}, и утверждение подтвердилось. Музыка была настоящей. Исполнения были убедительными. Педагогика была преобразующей.

Но PLOrk был построен в Принстоне. Каждая из пятнадцати исполнительских станций требовала MacBook (\$1.300+), заказного полусферического громкоговорителя, спроектированного Джеффом Снайдером~\citep{snyder2015hemispherical}, чтобы проецировать звук всенаправленно, как акустический инструмент, 8-канального аудиоусилителя, кондиционера питания, корпуса-стойки и набора контроллеров --- MIDI-клавиатур, графических планшетов, интерфейсов Wiimote, датчиков давления и наклона. Программный стек включал ChucK~\citep{wang2003chuck}, Max/MSP (по лицензии) и заказной сетевой код. Полная стоимость одного места: консервативно \$1.500, вероятно выше, если учесть труд по изготовлению и институциональные накладные расходы.

Для Принстона это не было барьером. Это было барьером для всех остальных.

В течение следующего десятилетия оркестры ноутбуков распространились по университетам --- SLOrk в Стэнфорде~\citep{wang2009designing}, L2Ork в Virginia Tech~\citep{bukvic2010l2ork}, CLOrk в Конкордии, BLOrk в Колорадо, HELO в Хаддерсфилде, OLO в Осло, SCLOrk в Санта-Кларе, CMU Laptop Orchestra в Карнеги-Меллон, ансамбли в Индиане, Georgia Tech, UBC и Национальном университете музыки в Перу. На пике, возможно, существовало около тридцати ансамблей по всему миру. Каждый жил внутри университета. Каждый требовал институционального финансирования, выделенного пространства, технического персонала и зачисленных студентов.

Движение оркестров ноутбуков достигло пика примерно в 500--1.000 участников в год по всему миру. Для сравнения: 1,2 миллиона американских детей играют в школьных оркестрах~\citep{warschauer2004technology}. Оркестр ноутбуков был прекрасной идеей, которая почти ни до кого не дошла.

Эта статья утверждает, что ограничивающим фактором никогда не было искусство --- PLOrk решил этот вопрос в первый же год --- а экономика и инфраструктура. Мы представляем \acos{}, операционную систему «на голом железе», которая превращает любой списанный ноутбук x86\_64 в полноценный музыкальный инструмент при стоимости на два порядка ниже модели PLOrk, и предлагаем, что возникающая в результате экономика делает \emph{планетарный} оркестр ноутбуков не только возможным, но --- учитывая объём доступного списанного железа --- почти неизбежным.

% ============ 2. ЛАНДШАФТ ОРКЕСТРОВ НОУТБУКОВ ============

\section{Ландшафт оркестров ноутбуков}
\label{sec:landscape}

\subsection{PLOrk: золотой стандарт}

Дизайн PLOrk был дотошным. Труман~\citep{trueman2007why} сформулировал центральную проблему: встроенные динамики ноутбука производят тонкий, бесформенный звук, который не может «заполнить зал так, как это делает акустический инструмент». Полусферический громкоговоритель --- геодезическая структура с шестью излучателями, направленными наружу по полусфере, --- был решением. Он дал каждой станции-ноутбуку физическое акустическое присутствие, \emph{тело} для её звука. Ноутбук стал, в рамках Трумана, не контроллером, запускающим сэмплы, а инструментом с пространственными характеристиками.

Это было музыкально верно и экономически катастрофично. Каждый полусферический громкоговоритель собирался вручную. Усилители, стойки и кондиционеры питания добавляли объём и стоимость. Результатом был ансамбль, который звучал необычайно и стоил целое состояние. Полностью оснащённая конфигурация PLOrk из 15 станций представляла собой вложение свыше \$30.000 только в оборудование --- до учёта зала, сетевой инфраструктуры, лицензий на ПО и технического персонала для поддержания всего этого.

\subsection{L2Ork: Linux-контрпредложение}

Ивица Ико Буквич в Virginia Tech распознал проблему доступа и построил L2Ork~\citep{bukvic2010l2ork}, Linux Laptop Orchestra, как явное снижение стоимости. L2Ork заменил MacBook на Linux-нетбуки, Max/MSP --- на Pd-L2Ork (заказной форк Pure Data), а принстонские полусферические громкоговорители --- на более дешёвые полусферические конструкции на стандартных излучателях. Результат: примерно \$750 за место, снижение на 50\%.

L2Ork был значимым вкладом. Настойчивость Буквича в отношении открытого ПО и доступного железа была философски согласована с миссией доступа. Но \$750 за место по-прежнему требуют институциональной закупки. Полусферические громкоговорители, даже более дешёвые, по-прежнему требуют изготовления. Pd-L2Ork по-прежнему требует установки, настройки и обслуживания на хостовом дистрибутиве Linux. L2Ork демократизировал оркестр ноутбуков внутри университетской системы. Он не вырвался из неё.

\subsection{Университетская ловушка}

Этот паттерн постоянен во всех тридцати с лишним оркестрах ноутбуков:

\begin{itemize}
  \item \textbf{Ограничен зачислением}: участвуют только зачисленные студенты
  \item \textbf{Финансируется кафедрами}: бюджеты на оборудование музыкальных или ИТ-кафедр
  \item \textbf{Зависит от пространства}: специальные репетиционные залы с сетевой инфраструктурой
  \item \textbf{Зависит от персонала}: преподаватели-руководители, технические ассистенты, ИТ-поддержка
  \item \textbf{Привязан к семестру}: ансамбли формируются и распускаются по академическому календарю
  \item \textbf{Географически фиксирован}: нужно физически находиться в университете
\end{itemize}

Это не случайные ограничения. Это структурные следствия модели затрат. Когда каждое место стоит \$750--1.500 и требует заказного железа, ансамбль \emph{должен} управляться институционально. Оборудование слишком дорогое, чтобы его терять, слишком специализированное, чтобы его заменять, и слишком сложное, чтобы обслуживать его без технического персонала. Университет --- не просто хозяин, он единственный жизнеспособный хозяин.

\subsection{Что говорят цифры}

\begin{table}[h]
\small
\centering
\begin{tabular}{lrrr}
\toprule
\textbf{Модель ансамбля} & \textbf{За место} & \textbf{30 мест} & \textbf{1000 мест} \\
\midrule
PLOrk (Princeton)       & \$1.500+ & \$45.000+ & невозможно \\
L2Ork (Virginia Tech)   & \$750      & \$22.500  & невозможно \\
SLOrk (Stanford)        & \$1.200+ & \$36.000+ & невозможно \\
Sonic Pi + Pi-железо    & \$150--200 & \$4.500   & \$150.000 \\
\textbf{AC OS + излишки} & \textbf{\$50--80} & \textbf{\$1.500} & \textbf{\$50.000} \\
\bottomrule
\end{tabular}
\caption{Сравнение затрат между моделями оркестра ноутбуков. «Невозможно» означает, что зависимость модели от заказного изготовления и институциональной инфраструктуры не даёт масштабироваться дальше $\sim$50 мест.}
\label{tab:cost}
\end{table}

Таблица рассказывает всю историю. По ценам PLOrk оркестр ноутбуков на 1.000 человек --- фантазия. По ценам AC~OS он стоит меньше одной конфигурации PLOrk. Разница не инкрементальна. Это фазовый переход.

% ============ 3. ЦУНАМИ ИЗЛИШКОВ ============

\section{Цунами излишков}
\label{sec:surplus}

\subsection{240 миллионов машин}

14 октября 2025 года Microsoft прекратила поддержку Windows~10~\citep{canalys2024windows}. Canalys оценила, что 240 миллионов ПК по всему миру --- машин, чьё железо полностью исправно, но чьи процессоры не удовлетворяют требованиям TPM~2.0 и поколению CPU для Windows~11 --- были объявлены «устаревшими» в одночасье. Это не старые машины. Это ноутбуки эпохи 2017--2020 годов с четырёхъядерными процессорами, 8\,ГБ ОЗУ, экранами 1080p и WiFi. Они стали электронными отходами не потому, что сломались, а потому, что Microsoft решила, что они закончили.

Это крупнейшее единичное событие появления излишков в истории персональных компьютеров. И оно приходит поверх уже существующих потоков излишков: возвраты корпоративного лизинга (циклы обновления 3--5 лет), ротация школьных Chromebook~\citep{pirg2023chromebook}, избыток розничных запасов и постоянный поток косметически повреждённых, но рабочих машин в каналы восстановления.

\subsection{Глобальный контекст электронных отходов}

Глобальный монитор электронных отходов ООН 2024~\citep{forti2024ewaste} сообщает о 62 миллионах тонн электронных отходов, образующихся в мире ежегодно, с приростом 2,6 миллиона тонн в год. Лишь 22,3\% задокументированы как правильно собранные и переработанные. Остальное оказывается на свалках, сжигается или экспортируется на неформальную переработку в страны Глобального Юга, где работники --- зачастую дети --- извлекают металлы в опасных условиях.

Каждый списанный ноутбук, отведённый из этого потока в музыкальный инструмент, --- небольшое экологическое вмешательство. В масштабе программа, повторно использующая всего лишь 0,1\% годового излишка ПК, производила бы 240.000 инструментов в год.

\subsection{Кладбище Chromebook}

Отчёт «Chromebook Churn» от U.S. PIRG Education Fund~\citep{pirg2023chromebook} документирует, как школьные округа по всей стране меняют Chromebook каждые 3--5 лет, тратя оценочно \$1,8 млрд на устройства, спроектированные под замену. Политика Auto Update Expiration (AUE) от Google лишает Chromebook поддержки после фиксированной даты независимо от состояния железа. Один только Oakland Unified управляет 40.000 Chromebook в этом цикле.

Ироничная трагедия: школы выбрасывают Chromebook, чьё железо превосходит требования к полноценному инструменту творческих вычислений, а затем тратят \$200--300 на новые. Округ, который перепрошил бы списанные Chromebook на \acos{}, мог бы превратить свой бюджет на электронные отходы в программу инструментов.

(Примечание: \acos{} в настоящее время ориентирован на машины x86\_64 UEFI. Большинство образовательных Chromebook используют процессоры ARM. Поддержка ARM --- первоочередная цель разработки, см.\ раздел~\ref{sec:future}. Однако Chromebook на Intel совместимы немедленно.)

\subsection{Цены на излишки}

Текущие цены (март 2026) на списанные ноутбуки на оптовых рынках восстановления и вторичных рынках:

\begin{table}[h]
\small
\centering
\begin{tabular}{lr}
\toprule
\textbf{Машина} & \textbf{Диапазон цен} \\
\midrule
ThinkPad X1 Carbon (6--8 поколение) & \$40--80 \\
ThinkPad T480/T490 & \$35--70 \\
HP EliteBook 840 G5/G6 & \$40--75 \\
Dell Latitude 7400/7410 & \$45--80 \\
Dell Latitude 5400 & \$30--55 \\
HP ProBook 440 G6/G7 & \$30--60 \\
Lenovo IdeaPad (разные) & \$25--50 \\
\bottomrule
\end{tabular}
\caption{Цены на списанные ноутбуки, март 2026. Цены отражают оптовые закупки (\textgreater 10 единиц) у сертифицированных восстановителей. Все модели включают экран 1080p, 4+ ядра, 8+\,ГБ ОЗУ, WiFi, полную клавиатуру и 4+ часа батареи.}
\label{tab:surplus}
\end{table}

Каждая из этих машин загружает \acos{} менее чем за 8 секунд. У каждой есть полная клавиатура, тачпад, микрофон, динамики, экран и WiFi. Каждая --- полноценный автономный музыкальный инструмент, ждущий файла на 89\,МБ.

% ============ 4. AC NATIVE OS КАК ИНСТРУМЕНТ ============

\section{AC Native OS как инструмент}
\label{sec:instrument}

Там, где PLOrk требовал собрать инструмент из компонентов --- ноутбук плюс громкоговоритель плюс усилитель плюс стойка плюс контроллеры плюс ПО --- \acos{} предлагает, что списанный ноутбук \emph{уже является} инструментом. Ему нужно лишь об этом сообщить.

\subsection{Архитектура вкратце}

\acos{} --- единый бинарник UEFI на 89\,МБ, содержащий ядро Linux, сжатый LZ4 initramfs и рантайм на C из 11.943 строк, работающий как PID~1~\citep{scudder2026os}. Нет пользовательского пространства, нет среды рабочего стола, нет менеджера пакетов, нет оболочки. Рантайм напрямую управляет экраном (DRM/KMS), звуком (ALSA на 192\,кГц), вводом (evdev), WiFi (wpa\_supplicant), камерой (V4L2) и речью (flite). Интерактивные программы («pieces») --- ES-модули JavaScript, исполняемые QuickJS~\citep{bellard2019quickjs}.

Время загрузки: \textbf{7,3 секунды} от кнопки питания до готового к игре инструмента.

\subsection{Встроенный синтез}

Рантайм включает 32-голосный полифонический синтезатор с пофолосными ADSR-огибающими, несколькими формами волны (синус, треугольник, пила, импульс, шум) и глобальным процессором ревербации. Звук работает на 192\,кГц через ALSA PCM --- тот же интерфейс, что используют профессиональные аудиостанции, без накладных расходов PulseAudio, PipeWire или JACK.

Это критическое расхождение с PLOrk. Полусферический громкоговоритель Трумана решал проблему пространственной проекции звука для ансамблей ноутбуков. \acos{} обходит проблему полностью: встроенные динамики ноутбука и есть динамики. Они не всенаправленные, не полусферические, не акустически идеальные. Но они \emph{есть} --- в каждом списанном ноутбуке, уже подключённые, не требующие изготовления, без усилителей, без кондиционеров питания, без стоек.

Для ансамбля PLOrk из 15 человек, играющего в концертном зале, полусферические громкоговорители --- правильный ответ. Для распределённого ансамбля из 1.000 человек, играющего в классах, общественных центрах, парках и гостиных, правильный ответ --- собственные динамики ноутбука. Полусферический громкоговоритель был решением проблемы концертного зала. Планетарный оркестр ноутбуков --- не концертный зал.

\subsection{У ноутбука уже есть всё}

\begin{table}[h]
\small
\centering
\begin{tabular}{ll}
\toprule
\textbf{PLOrk требует} & \textbf{У списанного ноутбука есть} \\
\midrule
MacBook & Любой ноутбук x86\_64 \\
Полусферический громкоговоритель & Встроенные динамики \\
8-канальный усилитель & Встроенный усилитель (ALSA) \\
Кондиционер питания & Адаптер AC или батарея \\
Корпус-стойка & Корпус ноутбука \\
MIDI-контроллер & Клавиатура + тачпад \\
Лицензия ChucK/Max/MSP & AC OS (бесплатно, встроено) \\
Сетевая инфраструктура & WiFi (встроен) \\
ИТ-персонал & Обновления OTA (автоматические) \\
\bottomrule
\end{tabular}
\caption{Требования к компонентам PLOrk против возможностей списанного ноутбука.}
\label{tab:plork-vs-surplus}
\end{table}

Списанный ноутбук --- не деградировавшая станция PLOrk. Это \emph{другой тип инструмента} --- автономный, портативный, питаемый батареей, в личной собственности и готовый к игре немедленно. То, чего ему не хватает в акустической проекции, он выигрывает в независимости. Станция PLOrk не может покинуть репетиционный зал. Ноутбук с AC~OS может пойти куда угодно.

\subsection{Сетевой ансамбль «из коробки»}

\ac{} включает встроенную многопользовательскую сеть: WebSocket для надёжных событий (вход, выход, счёт, состояние) и UDP для данных реального времени с низкой задержкой (позиция, скорость, непрерывные контроллеры). Любое устройство \acos{} в той же сети WiFi может присоединиться к сессии ансамбля без настройки. Сервер сессий~\citep{scudder2026ac} управляет маршрутизацией, и pieces могут определять произвольные сетевые протоколы для ансамблевого взаимодействия.

Сеть PLOrk требовала заказной инфраструктуры и технической настройки для каждого выступления. Ансамблевая сеть \acos{} всегда активна, встроена в слой ОС и работает в любом WiFi --- включая точку доступа на телефоне. Три человека со списанными ноутбуками в парке могут составить оркестр ноутбуков.

\subsection{Персонализация как идентичность инструмента}

Каждое устройство \acos{} загружается с отображённым и произнесённым вслух хэндлом владельца: «hi @maria». Цвета, токены и предпочтения хранятся в \texttt{config.json} на разделе EFI, сохраняясь через OTA-обновления~\citep{scudder2026os}. Машина --- не обобщённый терминал, это \emph{инструмент Марии}, и она это знает.

В ансамблевом контексте это означает, что у каждого инструмента есть имя. Дирижёр (или сеть) может обратиться к машине «@maria», направить к ней звук, назначить ей голос. Идентичность инструмента --- не дополнение; она фундаментальна. У каждого оркестрового инструмента есть имя. Теперь и у каждого ноутбука тоже.

% ============ 5. СРАВНЕНИЕ С ВЫДЕЛЕННЫМИ УСТРОЙСТВАМИ ============

\section{Ландшафт выделенных устройств}
\label{sec:devices}

Списанный ноутбук --- не единственное портативное вычислительное устройство в ценовом диапазоне \$50--300. Потребительский рынок уже принимает --- с энтузиазмом --- специально построенные вычислительные устройства, которые хорошо делают одно дело. Сравнение проясняет, чем \acos{} является и чем не является.

\subsection{Nintendo Switch (\$299 новая, \$150 б/у)}

Switch --- самое успешное выделенное устройство творческих вычислений в истории, хотя так и не позиционируется. Это портативный компьютер с сенсорным экраном, динамиками, WiFi, Bluetooth, датчиками движения, отсоединяемыми контроллерами и док-станцией для вывода на ТВ. Он работает на заказной ОС, которая загружается за секунды в единый интерфейс. Пользователи выбирают из библиотеки интерактивных программ («игр») по имени. Нет файлового менеджера, нет рабочего стола, нет терминала.

Структурное сходство с \acos{} поразительно: обе --- операционные системы единого назначения на портативном железе, обе быстро загружаются в интерактивный опыт, обе используют WiFi для многопользовательской игры. Switch продано 146 миллионов единиц~\citep{nintendo2024earnings} --- доказательство, что публика примет выделенные вычислительные устройства, когда опыт целостен, а барьер входа --- кнопка питания.

Разница в назначении и экономике. Switch стоит \$299 новой. Списанный ThinkPad с \acos{} стоит \$50. Switch --- закрытая платформа, принадлежащая Nintendo. \acos{} открыт, модифицируем и самообновляем. Switch не может быть перепрограммирован своим владельцем. Устройство \acos{} исполняет pieces, созданные пользователем. Но Switch доказывает форм-фактор: портативное, специально построенное, мгновенно включающееся вычислительное устройство --- это то, что люди хотят, понимают и используют ежедневно.

\subsection{Steam Deck (\$399 новая, \$250 б/у)}

Steam Deck от Valve --- портативный ПК x86\_64, работающий на заказном дистрибутиве Linux (SteamOS), который загружается прямо в библиотеку игр. Архитектурно это почти точно то, что делает \acos{}: специально построенная среда Linux на стандартном железе x86\_64, оптимизированная под единственный сценарий использования, с быстрой загрузкой и курируемым интерфейсом.

Steam Deck доказывает, что подход «Linux на голом железе x86\_64» работает коммерчески. Valve продано оценочно 5--10 миллионов единиц. Пользователи принимают и предпочитают выделенный опыт обобщённому рабочему столу. Steam Deck также доказывает, что встроенные динамики, хоть и акустически скромные, приемлемы для интерактивных медиа, когда опыт в остальном убедителен.

Там, где Steam Deck требует специально построенного железа (\$399), \acos{} работает на железе, которое уже существует и стоит \$50. Steam Deck исполняет игры; \acos{} исполняет pieces. Но концептуальная архитектура --- заказной Linux, быстрая загрузка, курируемая библиотека, встроенная сеть --- общая.

\subsection{Teenage Engineering (OP-1, EP-133, TX-6)}

Линейка продуктов Teenage Engineering представляет верхний предел выделенных творческих инструментов: OP-1 Field (\$1.999), EP-133 K.O. II (\$299), микшер TX-6 (\$1.199). Это специально построенные устройства для создания музыки с заказным железом, встроенными динамиками, батареей и интерфейсами с характером. Это красивые объекты. Они также, функционально, предметы роскоши. Основной рынок OP-1 --- не профессиональные музыканты и не студенты, а состоятельные взрослые-любители. Устройство чаще появляется в распаковках на YouTube и на стендах выставки NAMM, чем в классах или общественных центрах. За \$2.000 оно недоступно тем группам, которым творческие инструменты нужнее всего.

Это и есть паттерн: выделенное музыкальное железо ориентировано либо на институции (PLOrk, Ableton Push в университетских студиях), либо на состоятельных частных лиц (Teenage Engineering, Elektron, Moog). Середина --- доступные, программируемые, лично принадлежащие инструменты для людей без институциональной аффилиации и без свободных средств --- пуста.

\acos{} на списанном ноутбуке не может сравниться с промышленным дизайном Teenage Engineering. Клавиатура ThinkPad --- не секвенсорная сетка. Форм-фактор ноутбука --- не карманный инструмент. Но \acos{} предлагает то, что не может ни один роскошный инструмент: масштаб, программируемость и нулевую предельную стоимость. OP-1 для каждого ребёнка в школе стоит \$600.000 за 300 единиц. Списанный ноутбук для каждого ребёнка стоит \$15.000. OP-1 не может быть перепрограммирован владельцем. Инструмент \acos{} исполняет любой piece, который его владелец напишет или скачает. Вопрос не в том, так ли красив ThinkPad, как OP-1. Вопрос в том, должна ли красота быть предусловием доступа.

\subsection{Перспектива выделенных устройств}

Рынок потребительской электроники уже ответил на вопрос «Будут ли люди пользоваться выделенными вычислительными устройствами?» 146 миллионами Switch, 5--10 миллионами Steam Deck и процветающим рынком специально построенного музыкального железа. \acos{} применяет модель выделенного устройства к списанному железу: тот же архитектурный паттерн (быстрая загрузка, единое назначение, курируемый интерфейс, встроенная сеть), радикально другая экономика (бесплатное ПО на железе за \$50 вместо специально построенных устройств за \$300--400).

Списанный ноутбук --- не деградировавший компьютер общего назначения. Это \emph{нереализованное выделенное устройство} --- то, что всегда было способно быть инструментом, но вместо этого работало на Windows.

% ============ 6. АРГУМЕНТ МАСШТАБИРОВАНИЯ ============

\section{От ансамбля к планете}
\label{sec:scaling}

\subsection{Предел масштабирования PLOrk}

PLOrk обслуживает 15 человек в семестр. За двадцать лет, при полном зачислении каждый семестр, PLOrk обслужил примерно 600 отдельных участников. Это щедрая оценка --- многие студенты участвуют в нескольких семестрах, а ансамбль работает не в каждый период.

Масштабирование PLOrk до 100 мест потребовало бы \$150.000+ в оборудовании, большего исполнительского пространства, больше усилителей, больше электрической инфраструктуры, больше технического персонала. Масштабирование до 1.000 мест физически невозможно --- одни только полусферические громкоговорители потребовали бы склада.

Модель PLOrk масштабируется \emph{линейно по стоимости и логарифмически по сложности}. Каждое дополнительное место требует не только оборудования, но и пропорционально больше логистических накладных расходов. Это и есть фундаментальное ограничение: PLOrk --- ремесленный инструмент. Он прекрасен, но не может быть воспроизведён.

\subsection{Кривая масштабирования AC OS}

\acos{} масштабируется линейно по стоимости и \emph{сублинейно по сложности}. Каждый дополнительный инструмент --- ещё один списанный ноутбук и одна запись на USB --- шестьдесят секунд работы. Нет громкоговорителей, которые надо строить, нет стоек, которые надо кабелировать, нет ПО, которое надо устанавливать, нет учётных записей, которые надо заводить. Система OTA-обновлений означает, что каждый инструмент обслуживает себя сам. Добавление 1.001-го инструмента к флоту из 1.000 требует тех же усилий, что и добавление второго к флоту из одного.

\begin{table}[h]
\small
\centering
\begin{tabular}{lrr}
\toprule
\textbf{Масштаб} & \textbf{Модель PLOrk} & \textbf{Модель AC OS} \\
\midrule
15 мест & \$22.500 & \$750 \\
30 мест & \$45.000 & \$1.500 \\
100 мест & \$150.000 & \$5.000 \\
1.000 мест & невозможно & \$50.000 \\
10.000 мест & невозможно & \$500.000 \\
100.000 мест & невозможно & \$5.000.000 \\
\bottomrule
\end{tabular}
\caption{Затраты на масштабирование. Модель PLOrk становится физически невозможной свыше $\sim$50 мест из-за ограничений изготовления и пространства. Модель AC~OS масштабируется линейно только по стоимости железа; стоимость ПО и обслуживания равна нулю при любом масштабе.}
\label{tab:scaling}
\end{table}

При \$50 за инструмент общегородской оркестр ноутбуков из 1.000 инструментов стоит \$50.000 --- цена одной хорошо оснащённой установки PLOrk. Национальная программа, распространяющая 100.000 инструментов, стоит \$5 миллионов --- меньше, чем одно развёртывание OLPC на страну, которое обычно стоило \$10--50 миллионов~\citep{kraemer2009olpc} в железе, часто остававшемся неиспользованным.

\subsection{Оркестр на 240 миллионов мест}

Излишек от окончания жизни Windows~10 представляет 240 миллионов потенциальных инструментов. Не все они --- ноутбуки x86\_64. Не у всех рабочие батареи. Не все доступны. Но даже доля этого излишка, должным образом отведённая, затмевает любое развёртывание образовательных технологий в истории.

Рассмотрим: если бы 1\% списанных машин Windows~10 был перепрошит на \acos{}, результатом стали бы 2,4 миллиона творческих инструментов, распределённых по всему миру. Проект OLPC на пике развернул примерно 3 миллиона ноутбуков XO за весь срок службы --- ценой в миллиарды. Один процент излишка Windows~10, при нулевой стоимости изготовления и \$0 за единицу ПО, сравнивается с полным развёртыванием OLPC.

Инструменты уже существуют. Они в шкафах, на складах и в контейнерах для переработки. ПО для их активации --- 89\,МБ.

% ============ 6. PLORK ПЛАНЕТУ ============

\section{PLOrk'ить планету: конкретное предложение}
\label{sec:proposal}

\subsection{Общественные события перепрошивки}

Простейшая модель развёртывания: общественное событие, где участники приносят списанные ноутбуки и уходят с инструментами. Событие требует одного человека с USB-накопителями, шаблона \texttt{config.json} и стола. Каждая машина персонализируется («Как тебя зовут?»), перепрошивается (шестьдесят секунд), загружается (семь секунд) и возвращается. Участник уходит с личным творческим инструментом, который будет самообновляться вечно.

Это не умозрительно. Это простая логистическая операция --- проще кампании регистрации избирателей, дешевле школьной благотворительной ярмарки выпечки.

\subsection{Школьные программы без бюджета на оборудование}

Школьный округ, списывающий Chromebook или машины Windows~10, может перенаправить железо в программу инструментов вместо переработчика. Chromebook на базе Intel и все ПК x86\_64 совместимы немедленно. «Бюджет на оборудование» равен нулю --- оборудование уже существует и уже оплачено. Единственная стоимость --- труд по перепрошивке USB-накопителей.

Для школ, уже выбросивших своё железо: списанные ThinkPad от корпоративных доноров по \$30--50 каждый дают классный ансамбль из 30 инструментов за \$900--1.500. Это меньше годовой подписки на ПО для большинства школьных технологических программ.

\subsection{Внеклассные и общественные ансамбли}

Оркестр ноутбуков никогда не существовал вне университета, потому что никогда не был экономически жизнеспособен вне университета. При \$50 за инструмент без повторяющихся расходов ансамбли ноутбуков становятся жизнеспособны для:

\begin{itemize}
  \item Внеклассных музыкальных программ
  \item Общественных центров и библиотек
  \item Молодёжных оркестров и мейкерспейсов
  \item Центров для пожилых и домов престарелых
  \item Заключённых (для локального ансамбля интернет не требуется)
  \item Лагерей беженцев и помощи при бедствиях (работает от батареи, WiFi опционален)
  \item Сельских и удалённых сообществ (OTA-обновления через любое соединение)
\end{itemize}

Ни один из этих контекстов не смог бы поддержать станцию PLOrk. Все могут поддержать списанный ноутбук.

\subsection{Распределённый оркестр}

Поскольку ансамблевые сети \acos{} работают поверх стандартного WiFi и интернет-соединений, оркестру ноутбуков не нужно быть в одной комнате. Piece, спроектированный для распределённого исполнения --- где задержка является композиционным параметром, а не дефектом --- позволяет:

\begin{itemize}
  \item Одновременные общегородские исполнения
  \item Межокружные ансамбли «школа--школа»
  \item Трансграничные телематические оркестровые события
  \item Асинхронные «эстафетные» композиции, передаваемые между часовыми поясами
\end{itemize}

Concordia Laptop Orchestra (CLOrk) экспериментировала с телематическим исполнением в академических контекстах. \acos{} делает телематическое исполнение доступным для любого, у кого есть WiFi и списанный ноутбук.

% ============ 7. ЧТО PLORK СДЕЛАЛ ХОРОШО ============

\section{Что PLOrk сделал хорошо (и что упустил)}
\label{sec:lessons}

Эта статья не критика художественного или педагогического качества PLOrk. PLOrk --- одно из важнейших достижений компьютерной музыки, и его вклад фундаментален:

\begin{enumerate}
  \item \textbf{Легитимность}: PLOrk доказал, что ансамбли ноутбуков производят настоящую музыку, а не трюки. В 2005 году это не было очевидно.
  \item \textbf{Пространственность}: Полусферический громкоговоритель показал, что электронные инструменты выигрывают от акустического присутствия. Это понимание сохраняется, даже когда конкретное решение меняется.
  \item \textbf{Композиция}: PLOrk породил новый репертуар~\citep{smallwood2008composing} и новые композиционные стратегии для сетевых инструментов. Этот корпус работ напрямую применим к любому ансамблю ноутбуков любого масштаба.
  \item \textbf{Педагогика}: Структура курса PLOrk --- сочетающая исполнение, композицию, дизайн инструментов и акустическое исследование --- модель для образования в творческих вычислениях~\citep{trueman2007why}.
  \item \textbf{Сообщество}: PLOrk создал новую социальную форму --- оркестр ноутбуков как ансамбль --- которая породила радость, сотрудничество и художественное открытие.
\end{enumerate}

Что PLOrk упустил --- и что, справедливости ради, не было его миссией решать --- это доступ. PLOrk был исследовательским проектом в элитном университете. Он был спроектирован, чтобы ответить на вопрос «Могут ли ансамбли ноутбуков быть музыкально значимыми?» Он не был спроектирован, чтобы ответить «Может ли каждый иметь такой?»

L2Ork продвинулся к вопросу доступа более дешёвым железом и открытым ПО, но структурная зависимость от институциональной инфраструктуры осталась. Sonic Pi~\citep{aaron2016sonic} решил вопрос индивидуального доступа к музыкальному кодингу, но это однопользовательский инструмент на хостовой ОС, а не ансамблевый инструмент. Ни один из этих проектов не предложил модель, где оркестр ноутбуков может работать вообще без какой-либо институции.

Это и есть предложение здесь. Не заменить PLOrk, а исполнить мечту, которую PLOrk начал: что сеть компьютеров может быть новым видом оркестра. PLOrk доказал, что может. \acos{} предлагает сделать это настолько дёшево, что доказательство станет планетарным.

% ============ 8. МУЗЫКА И ПОЛИТЭКОНОМИЯ ============

\section{Музыка, доступ и политэкономия}
\label{sec:theory}

Жак Аттали в книге \emph{Noise}~\citep{attali1985noise} утверждал, что музыка пророческа: изменения в производстве и распространении музыки предвосхищают изменения в политэкономии. История музыки движется от ритуала (жертвоприношение) к репрезентации (концертный зал), к повторению (запись) и к тому, что Аттали назвал \emph{композицией} --- режиму, где все творят, где различие между исполнителем и слушателем растворяется.

Оркестр ноутбуков, в рамках Аттали, --- ансамбль эпохи композиции. Нет аудитории в традиционном смысле; участники одновременно композиторы, исполнители и слушатели. Но если музыка эпохи композиции определяется всеобщим участием, то университетский оркестр ноутбуков --- противоречие: он ограничивает участие зачислением, платой за обучение и географией. Это музыка эпохи композиции с контролем доступа эпохи репрезентации.

Томас Турино~\citep{turino2008music} различает музыку \emph{партиципаторную} (где успех измеряется степенью участия) и музыку \emph{презентационную} (где успех измеряется качеством исполнения для аудитории). Выступления PLOrk презентационны --- они проходят на концертных сценах для сидящих слушателей. Но \emph{репетиция} PLOrk партиципаторна --- играют все, переживают все, коллективный звук и есть смысл.

Планетарный оркестр ноутбуков партиципаторен по умолчанию. Нет сцены. Нет аудитории. Есть только участники, у каждого свой инструмент, подключающиеся когда и как они выберут. Рамка Турино предсказывает, что этот сдвиг --- от презентационного к партиципаторному оркестру ноутбуков --- изменит не только то, кто играет, но и то, что «играть» означает.

«Инструменты для дружелюбия» Ивана Иллича~\citep{illich1973tools} завершают теоретическую рамку. Дружелюбный инструмент --- тот, который может быть использован кем угодно, без специальной подготовки, для достижения целей, которые они сами определяют. Станция PLOrk --- с её полусферическим громкоговорителем, стойкой усилителя и патчем ChucK --- не дружелюбный инструмент. Она требует подготовки, инфраструктуры и институциональной поддержки. Списанный ноутбук с \acos{} приближается к дружелюбию: он загружается в готовый к игре инструмент, не требует настройки, обновляется сам и принадлежит своему владельцу.

% ============ 9. СВЯЗАННЫЕ РАБОТЫ ============

\section{Связанные работы}
\label{sec:related}

\textbf{Оркестры ноутбуков.} PLOrk~\citep{trueman2006plork} установил форму. L2Ork~\citep{bukvic2010l2ork} доказал жизнеспособность Linux и снижение стоимости. SLOrk~\citep{wang2009designing} внёс ChucK-педагогику и громкоговоритель из салатницы IKEA. С 2005 года работало более тридцати университетских ансамблей. Ни один не работал вне университета.

\textbf{Музыкальный кодинг для доступа.} Sonic Pi~\citep{aaron2016sonic} работает на Raspberry Pi и настольных ОС, обучая музыке через live coding. Это ближайший существующий проект к демократизированной компьютерной музыке, но он однопользовательский, требует хостовой ОС и не является ансамблевым инструментом.

\textbf{Образовательное железо.} OLPC~\citep{negroponte1995being, kraemer2009olpc} продемонстрировал и амбицию, и режим провала железоцентричного доступа. Raspberry Pi~\citep{raspberry2012} снизил стоимость платы, но не полную стоимость системы. Chromebook~\citep{pirg2023chromebook} достигли школьного развёртывания, но с запланированным устареванием и зависимостью от поставщика.

\textbf{Платформы творческих вычислений.} Processing~\citep{reas2007processing} и Scratch~\citep{resnick2009scratch} снизили барьеры творческого кодинга, но предполагают ОС общего назначения. \ac{}~\citep{scudder2026ac} предоставляет веб-платформу; \acos{}~\citep{scudder2026os} устраняет зависимость от ОС полностью.

\textbf{Теория.} Аттали~\citep{attali1985noise} о музыке как политическом пророчестве; Турино~\citep{turino2008music} о партиципаторном против презентационного исполнения; Иллич~\citep{illich1973tools} о дружелюбных инструментах; Паперт~\citep{papert1980mindstorms} о конструкционистском обучении; Киттлер~\citep{kittler1995no} о материальности вычислений; Маклюэн~\citep{mcluhan1964understanding} о медиа как среде; Юкелес~\citep{ukeles1969manifesto} об обслуживании как творческой практике.

% ============ 10. OLPC: УРОКИ ПОСЛЕДНЕЙ ПЛАНЕТАРНОЙ АМБИЦИИ ============

\section{Уроки OLPC}
\label{sec:olpc}

Проект One Laptop Per Child~\citep{negroponte1995being, kraemer2009olpc} --- ближайший прецедент к описанной здесь амбиции --- и самый поучительный провал. Николас Негропонте объявил в 2005 году (в тот же год основания PLOrk), что MIT произведёт ноутбук за \$100 и распространит его каждому ребёнку в развивающемся мире. Ноутбук XO вышел в 2007 году по \$188. За срок службы OLPC развернул примерно 3 миллиона единиц примерно в 40 странах, при общей оценочной стоимости программы в миллиарды, когда учтены логистика, обучение и государственные партнёрства.

OLPC провалился по причинам, теперь хорошо задокументированным~\citep{kraemer2009olpc}: железо было специально построенным и незнакомым (новый форм-фактор, новая ОС, новый интерфейс); локальный ремонт был невозможен, потому что детали были проприетарными; ПО (Sugar OS) не имело экосистемы или преемственности с тем, что учителя или ученики уже знали; развёртывания требовали государственной поддержки, министерской координации и логистических цепочек, которые рассыпались в поле; и стоимость за единицу, хоть и ниже коммерческих ноутбуков, оставалась слишком высокой для беднейших контекстов при умножении на миллионы.

\acos{} обращает каждый из этих режимов провала:

\begin{table}[h]
\small
\centering
\begin{tabularx}{\columnwidth}{lXX}
\toprule
\textbf{Режим провала OLPC} & \textbf{OLPC} & \textbf{AC OS} \\
\midrule
Железо & Специальный XO (\$188) & Списанные ноутбуки (\$50), уже изготовлены \\
Форм-фактор & Незнаком пользователям & Стандартный ноутбук, уже знаком \\
Ремонт & Проприетарные детали & Стандартное железо, ремонтируется локально \\
ПО & Sugar OS (новое, без экосистемы) & AC OS (целевое, но самообновляемое) \\
Развёртывание & Государственно-министерская координация & Один человек с USB-накопителями \\
Стоимость за единицу & \$188 + логистика & \$50 железо + \$0 ПО \\
Достигнутый масштаб & 3 млн единиц за 10 лет & 240 млн машин уже существуют \\
\bottomrule
\end{tabularx}
\caption{Режимы провала OLPC и их соответствия в AC OS.}
\label{tab:olpc}
\end{table}

Самый глубокий урок OLPC в том, что железоцентричные развёртывания проваливаются, когда требуют изготовления, логистики и институциональной координации. \acos{} не требует ничего из этого, потому что железо уже существует, уже глобально распределено и уже выбрасывается. Механизм развёртывания --- файл, а не цепочка поставок.

% ============ 11. ОГРАНИЧЕНИЯ ============

\section{Ограничения}
\label{sec:limitations}

Эта статья выдвигает амбициозное утверждение. Ограничения реальны и должны быть изложены ясно.

\textbf{Качество звука.} Встроенные динамики ноутбуков акустически слабы. У них нет басового отклика, ограниченный динамический диапазон, и они проецируют звук узко. Полусферический громкоговоритель Трумана решал реальную проблему --- электронным инструментам нужно акустическое присутствие, чтобы работать музыкально в ансамблевых контекстах. Динамики списанного ThinkPad не могут сравниться с этим. Для интимных или малогрупповых контекстов встроенных динамиков может быть достаточно. Для любого ансамбля больше комнаты нужна внешняя амплификация, что частично подрывает утверждение «без дополнительного железа».

\textbf{Поддержка ARM.} \acos{} в настоящее время ориентирован только на машины x86\_64 UEFI. Большинство списанных образовательных Chromebook используют процессоры ARM (Mediatek, Qualcomm). Пока поддержка ARM недоступна, крупнейший единый массив списанного образовательного железа остаётся недостижимым. Это самый критический инженерный пробел.

\textbf{Без формальной оценки.} Эта статья не представляет ни пользовательского исследования, ни контролируемого эксперимента, ни систематической оценки музыкального качества, результатов обучения или ансамблевой координации. Утверждения о доступе, стоимости и масштабируемости основаны на экономике железа и логистическом анализе, а не на эмпирическом педагогическом исследовании. Строгая оценка --- сравнивающая музыкальные результаты ансамбля типа PLOrk и ансамбля \acos{} эквивалентного размера --- существенно усилила бы аргумент.

\textbf{Проект одного автора.} \acos{} в основном написан одним человеком. Bus-фактор равен единице. Планетарное развёртывание не может зависеть от единственного сопровождающего. Система должна быть задокументирована, кодовая база должна быть читаема для новых контрибьюторов, а инфраструктура OTA-обновлений должна быть устойчива к отсутствию её создателя.

\textbf{Сетевая задержка.} Распределённое ансамблевое исполнение через интернет вносит задержку, которая фундаментально меняет, какая музыка возможна. Локальная задержка WiFi (1--5\,мс) позволяет плотную синхронизацию. Интернет-задержка (50--200\,мс) --- нет. Статья указывает на «задержку как композиционный параметр», но не демонстрирует композиций, которые заставляют это работать музыкально. Это открытая исследовательская проблема, не решённая.

\textbf{Без учебной программы.} Педагогика PLOrk~\citep{trueman2007why} разрабатывалась годами преподавания. У \acos{} нет учебной программы, нет планов уроков, нет задокументированной практики преподавания. Инструмент существует; педагогика нет. Без неё развёртывание железа в школах рискует повторить паттерн OLPC --- раздача устройств без образовательной инфраструктуры, делающей их значимыми.

% ============ 12. БУДУЩАЯ РАБОТА ============

\section{Следствия и будущая работа}
\label{sec:future}

\subsection{Поддержка ARM для Chromebook}

Инженерная цель наибольшего влияния --- поддержка ARM для Chromebook на Mediatek и Qualcomm --- самого распространённого списанного образовательного железа в мире. Достижение этого разблокировало бы миллионы дополнительных инструментов из одного только школьного излишка.

\subsection{Встроенные инструменты композиции}

Текущие pieces \acos{} включают хроматическую клавиатуру, навигатор команд и системную панель. Ансамблевый сценарий требует секвенсоров, редакторов паттернов, инструментов лупинга и интерфейсов совместной композиции --- всё как pieces, заменяемые во время исполнения.

\subsection{Пространственность без громкоговорителей}

Понимание PLOrk о пространственном звуке остаётся верным, даже когда решение меняется. Будущая работа включает программно определяемую пространственность: использование распределённых физических позиций ансамбля (сообщаемых через UDP) для создания эмерджентных полей пространственного звука из совокупности множества динамиков ноутбуков. Кольцо из 30 ноутбуков \emph{есть} полусферический громкоговоритель, если ПО знает, где каждый из них.

\subsection{Разработка учебной программы}

Педагогика PLOrk~\citep{trueman2007why} --- сочетающая исполнение, композицию, дизайн инструментов и акустическое исследование --- должна быть адаптирована к контекстам нулевой инфраструктуры. Это значит учебные программы, не предполагающие выделенного пространства, технического персонала, предшествующего опыта программирования, и инструменты, которыми участники владеют и забирают домой.

\subsection{Mesh-сети}

Одноранговое ансамблевое соединение без интернет-инфраструктуры позволило бы выступления в местах без WiFi: на открытом воздухе, в сельской местности, в развивающихся регионах. Bluetooth-mesh или ad-hoc WiFi между устройствами \acos{} --- активная цель разработки.

% ============ 13. ЗАКЛЮЧЕНИЕ ============

\section{Заключение}

В 2005 году Дэн Труман спросил «Зачем оркестр ноутбуков?»~\citep{trueman2007why} и дал убедительный ответ: потому что сетевые компьютеры могут быть новым видом ансамбля, с новыми тембрами, новыми пространственными характеристиками, новыми композиционными возможностями и новыми социальными конфигурациями. Он был прав. PLOrk это доказал.

Двадцать лет спустя вопрос другой. Не «зачем оркестр ноутбуков?», а «почему только \emph{один}?» Почему тридцать, в тридцати университетах, обслуживающие несколько сотен человек в год? Почему не триста, или три тысячи, или три миллиона?

Ответ был экономическим. Полусферические громкоговорители стоят денег. Стойки усилителей стоят денег. MacBook стоят денег. Лицензии Max/MSP стоят денег. Технический персонал стоит денег. Оркестр ноутбуков был дорогим, поэтому он остался в единственном месте, где дорогие вещи для немногих терпят: в университете.

\acos{} устраняет экономический барьер. Файл на 89\,МБ. Списанный ноутбук за \$50. Семь секунд до загрузки. Встроенный синтез, встроенная сеть, встроенная персонализация, встроенные обновления. Никаких громкоговорителей, которые надо строить, никакого ПО, которое надо устанавливать, никакого ИТ-отдела, который надо содержать, никакой институции, которая должна это размещать.

С 240 миллионами ПК, входящих на рынок излишков только из-за Windows~10, и 62 миллионами тонн электронных отходов, накапливающихся ежегодно~\citep{forti2024ewaste}, сырьё для планетарного оркестра ноутбуков не дефицитно --- оно постыдно изобильно. Железо существует. ПО существует. Сеть существует. Музыка, как PLOrk показал нам двадцать лет назад, реальна.

«PLOrk'ить планету» --- не метафора. Это логистическая задача. И логистика только что подешевела в сто раз.

\vspace{0.5em}
\noindent\textbf{ORCID:} \href{https://orcid.org/0009-0007-4460-4913}{0009-0007-4460-4913}

\vspace{0.5em}
\noindent\textit{Перевод с английского. Оригинальная версия доступна на \url{https://papers.aesthetic.computer}}

% ============ ССЫЛКИ ============

\bibliographystyle{plainnat}
\bibliography{references}

\end{document}
